\begin{definition}
    Процесс $(X_t, t \in T)$~--- \textit{модификация} процесса $(Y_t, t \in T)$, если $\forall t \in T \  P(X_t = Y_t) = 1$.  При этом эти процессы называются стохастически эквивалентными.
\end{definition}


\begin{theorem}[о непрерывной модификации, Колмогоров]

Пусть процесс $(X_t, t \in [a, b])$ таков, что $\exists c, \eps, \alpha > 0$ такие, что
$$
\forall t, s \in [a, b] \quad \E |X_t - X_s|^\alpha \leq c \cdot |t - s|^{1 + \varepsilon}
$$

Тогда у процесса $X_t$ существует стохастически эквивалентный процесс, почти все
траектории которого непрерывны.
\end{theorem}


\begin{statement}\label{8.3}
     У винеровского процесса существует непрерывная модификация (модификация с непрерывными траекториями).
\end{statement}

\begin{proof}
Пусть $\alpha = 4$, тогда $\E (W_t - W_s)^4 = 3(t - s)^2$. Тогда по теореме Колмогорова $W_t$ допускает непрерывную модификацию на любом конечном отрезке. Пусть $(W_t^{(n)}, t \in [n, n + 1])$~--- непрерывная модификация $W_t$ на $[n, n + 1]$, $n \in \Z_+$.

Рассмотрим $\widetilde{W}_t =\{W_t^{(n)}, \ t \in [n, n + 1), \  n \in \Z_+\}$. Полученная модификация может быть разрывна только в целых точках (в остальных она строго непрерывна), когда $W_{n+1}^{(n)} \neq W_{n+1}^{(n + 1)}$. Но
$$
\begin{cases}
    P\brackets{W_{n + 1}^{(n)} = W_{n + 1}} = 1 \\
    P\brackets{W_{n + 1}^{(n + 1)} = W_{n + 1}} = 1 \\
\end{cases}
\Longrightarrow P\brackets{W_{n + 1}^{(n)} = W_{n + 1}^{(n + 1)}} = 1 \Longrightarrow P\brackets{\exists n: \ W_{n + 1}^{(n)} \ne W_{n + 1}^{(n + 1)}} = 0
$$
Следовательно, $\widetilde{W}_t$ имеет п.н. непрерывные траектории. Тогда положим
$$
\widetilde{\widetilde{W}}_t (\omega) = 
\begin{cases}
    \widetilde{W}_t (\omega), & \text{если траектория $\widetilde{W}_t(\omega)$ непрерывна} \\
    0, & \text{иначе}
\end{cases}
$$

Тогда $\widetilde{\widetilde{W}}_t$ является модификацией $\widetilde{W}_t$, и, как следствие, является модификацией $W_t$.\\
Значит, построили нужную непрерывную модификацию.
\end{proof} 

\begin{theorem}
    Пусть $\{e_n\}$ — ортонормированный базис в вещественном пространстве $L^2[0, T]$ и $\{\xi_n\}$ — последовательность независимых стандартных гауссовских случайных величин. Тогда ряд 
    \[ \sum_{n=1}^\infty \xi_n(\omega)\varphi_n(t), \quad \varphi_n(t) = \int_0^t e_n(s) ds \]
    
    равномерно сходится на $[0, T]$ при п.в. $\omega$, причем его сумма $w_t(\omega)$ есть винеровский процесс.
\end{theorem}